\honest{Replace the Symbol with the Substance}{Eliezer Yudkowsky}{2008}

Yesterday I described baseball without the word ``baseball.'' To do that, you visualize. A bat is a ``long, round, tapering, wooden rod''; a ball is a ``leather-covered spheroid''; the bases are fixed positions that players run to because the rules make them safe.

The obstacle is the word itself, which ``rushes in and obliterates the details you're trying to see.'' Playing Taboo is one of the ``fundamental'' rationalist capacities, on the same level as asking ``Why?'' I give no reason for the ranking.

To show how it keeps a purpose in view, I quote my own ``Lost Purposes'': a student studies ``12th-century knitting patterns'' because the degree requires it. Taboo ``school,'' ``education'' and ``degree,'' and you will notice that school ``currently seems to consist of sitting next to bored teenagers,'' that a degree is ``a piece of paper with some writing on it,'' and that education ``is forgetting the material as soon as you're tested on it.''

My new step is this. Schools that must produce test scores teach to the test, and a mental category that only partly matches your goal can fail the same way inside your head. If you want to learn and ``education'' is your name for it, you may not notice whether you are learning. \nb{The institutional half is Campbell's law; see the note on ``Lost Purposes.'' The internal half rests on this one example.}

To categorize is to throw away information. Even ``play Taboo'' is a leaky category: the board game bans only five words. The real version is played with a voluntary handicap. You also refuse unlisted synonyms and any new handle for the old idea. Judge the result by whether it serves its purpose. My instruction, in eleven phrasings, is to replace the symbol with the substance.

My example is ``The Simple Truth,'' written to describe truth ``without invoking terms like'' ``believe'' or ``real.'' The word ``true'' appears in it, I admit, but not to define truth. \nb{The published text uses forms of ``believe'' and ``belief'' 54 times and ``reality'' 15 times, including in the narrator's own account: ``The correspondence between reality and my beliefs comes from reality controlling my beliefs, not the other way around.'' It does also describe reality at a lower level, as ``whatever-it-is that determines my experimental results.''} Bayes's rule describes evidence in pure math; try to define such terms in words ``and you'll just go in circles.''

Last comes ``the most important word of all to Taboo,'' which I do not name. I have warned before against overusing it. Your understanding is measured by whether you can say what you are doing ``without using that word or any of its synonyms.''
